\documentclass[11pt,a4paper]{article}

\usepackage[utf8]{inputenc}
\usepackage[T1]{fontenc}
\usepackage[italian]{babel}
\usepackage[a4paper,margin=2.4cm]{geometry}
\usepackage{enumitem}
\usepackage{xcolor}
\usepackage{titlesec}
\usepackage[colorlinks=true,linkcolor=black,urlcolor=blue]{hyperref}

\definecolor{verde}{RGB}{46,125,80}
\titleformat{\section}{\Large\bfseries\color{verde}}{\thesection}{0.6em}{}
\titleformat{\subsection}{\large\bfseries}{\thesubsection}{0.5em}{}

\newcommand{\code}[1]{\texttt{#1}}

\title{\textbf{Relazione tecnica del progetto}\\[2pt]
\large Architettura, tecnologie e scelte progettuali}
\author{Gruppo Mu.To}
\date{A.A. 2025--26}

\begin{document}
\maketitle

\noindent
Questo documento descrive le \emph{tecnologie}, l'\emph{architettura} e la
\emph{struttura} del sito, motivando le scelte progettuali. Non descrive il
dominio applicativo: l'obiettivo è spiegare \emph{come} il sistema è costruito e
\emph{perché}, in funzione dei criteri di valutazione (generalità dei tool,
flessibilità, usabilità, sofisticazione grafica).

\tableofcontents
\bigskip

% =====================================================================
\section{Visione d'insieme dell'architettura}

Il sistema è organizzato a \textbf{livelli} (\emph{tiered}), con una netta
separazione tra dato, logica applicativa e presentazione. I componenti
comunicano esclusivamente via HTTP/JSON, così che ciascuno sia sostituibile in
modo indipendente.

\begin{verbatim}
   Browser
      |  HTTPS (terminato dal proxy del cluster)
      v
+-----------------------+   serve le 3 SPA compilate (dist/)
|  BFF  (Express)       |
|  server_bff.js :8000  |---- proxy /api ----+
+-----------------------+   (HTTPS interno)  |
      |  spawn interni                       v
      |                          +----------------------+   driver   +----------------+
      |                          |  openAPI (Express)   |  mongodb   |   MongoDB      |
      |                          |  REST API :3000      |----------->| (container DB) |
      v                          +----------+-----------+            +----------------+
+----------------------+                    |  ^
| SVG server :3001     |<-------------------+  |
| genera mappe SVG     |   HTTPS interno       v
+----------------------+              Provider IA (Groq, OpenAI-compatibile)
\end{verbatim}

\subsection{I componenti}
\begin{description}[leftmargin=1.4em,style=nextline]
  \item[openAPI (\code{server/openAPI})] È il cuore applicativo: una REST API in
    \textbf{Node.js + Express} che possiede la logica di dominio (musei, oggetti,
    percorsi, utenti, visite guidate, marketplace) e l'unico accesso a
    \textbf{MongoDB}. Espone inoltre l'integrazione con l'IA.
  \item[BFF (\code{navigator/UI/bff})] Il \emph{Backend-for-Frontend}: un secondo
    server Express che (a)~serve le tre Single Page Application già compilate e
    (b)~fa da \textbf{reverse proxy} verso l'API sul prefisso \code{/api},
    iniettando la chiave applicativa. È l'unico processo esposto pubblicamente.
  \item[SVG server (\code{map-creator/js\_server})] Servizio dedicato che, a
    partire dal layout di un museo, costruisce la \textbf{mappa vettoriale SVG}
    delle sale; è isolato perché la generazione grafica è una responsabilità
    distinta dal resto della logica.
  \item[Le tre SPA] \emph{marketplace} (catalogo, acquisti, visite),
    \emph{viewer/navigator} (la visita vera e propria sulla mappa) ed
    \emph{editor} (creazione/modifica dei musei).
\end{description}

Questa scomposizione è una scelta di \textbf{generalità}: ogni servizio fa una
cosa sola, dialoga via JSON e potrebbe essere riscritto in un altro linguaggio
senza toccare gli altri. Il BFF, in particolare, disaccoppia il browser
dall'API: il client non conosce né l'indirizzo reale dell'API né la sua chiave.

% =====================================================================
\section{Stack tecnologico e motivazione delle scelte (generalità)}

\subsection{Lato server}
\begin{itemize}[leftmargin=1.3em]
  \item \textbf{Node.js + Express} per API, BFF e SVG server: un unico runtime e
    un unico modello di routing su tutti i servizi server, che riduce il costo
    cognitivo e favorisce il riuso.
  \item \textbf{MongoDB} tramite il \emph{driver nativo} (\code{mongodb}), non un
    ORM. La scelta è deliberata: i documenti dei musei hanno una forma
    annidata e variabile (sale, oggetti, descrizioni multidimensionali) che mal
    si presta a uno schema rigido. Il modello a documenti rende naturale
    rappresentare un museo come un singolo aggregato.
  \item \textbf{Elaborazione media e QR}: \code{sharp} per ridimensionare e
    normalizzare le immagini caricate, \code{qrcode} per generare i codici QR di
    sblocco contenuti, \code{multer} per gli upload \emph{multipart}.
  \item \textbf{IA}: provider \textbf{Groq} con modello
    \code{llama-3.3-70b-versatile}, contattato tramite l'API
    \emph{OpenAI-compatibile}. Usare un'interfaccia standard (non un SDK
    proprietario) è una scelta di generalità: cambiare provider (OpenAI, Ollama
    locale, \dots) richiede solo variabili d'ambiente, non codice.
\end{itemize}

\subsection{Lato client: tre SPA, due filosofie}
La parte più significativa è la scelta differenziata del livello di
presentazione, calibrata sulla complessità di ciascuna applicazione.

\begin{itemize}[leftmargin=1.3em]
  \item \textbf{viewer (navigator)} ed \textbf{editor} usano \textbf{React}
    (il viewer in \textbf{TypeScript}). Sono applicazioni a stato ricco e
    fortemente interattive --- la navigazione sulla mappa SVG con \emph{pathfinding},
    la chat, l'editor di musei --- dove un framework a componenti con stato
    reattivo e tipizzazione statica ripaga la sua complessità.
  \item \textbf{marketplace} usa \textbf{Alpine.js} + JavaScript ``vanilla''.
    È stato \emph{deliberatamente riscritto} per non dipendere da un framework
    pesante: la reattività dichiarativa di Alpine (\code{x-data},
    \code{x-for}, \code{x-model}) copre esattamente i bisogni di una vetrina con
    filtri, modali e form, senza il peso di React. Il risultato è un bundle
    circa il 60\% più leggero. Questo dimostra \textbf{uso corretto della
    tecnologia minima sufficiente} invece di adottare un framework per inerzia.
  \item \textbf{Build}: tutte e tre con \textbf{Vite}, usato solo come
    bundler/dev-server (non è un framework UI). In sviluppo, il proxy di Vite
    inoltra \code{/api} all'API iniettando la chiave, replicando il
    comportamento del BFF in produzione.
\end{itemize}

% =====================================================================
\section{Organizzazione del codice e flessibilità}

\subsection{Struttura dei sorgenti}
Ogni applicazione vive in una propria directory con sorgenti in \code{src/} e
output di build in \code{dist/}. La separazione \code{src}/\code{dist} mantiene
puliti i sorgenti versionati; il BFF serve i \code{dist/} come contenuto
statico. La logica di rete (\code{api.js}) e la traduzione (\code{mpLocales.js})
sono moduli \emph{framework-agnostic}: durante la migrazione del marketplace da
React ad Alpine sono stati riusati invariati, segno di un buon disaccoppiamento.

\subsection{Modello dei dati e delle annotazioni}
Un oggetto museale non ha una singola descrizione, ma una \textbf{matrice di
descrizioni} indicizzata su due dimensioni:
\begin{itemize}[leftmargin=1.3em,nosep]
  \item \textbf{livello} del visitatore (bambino, studente, esperto, avanzato);
  \item \textbf{durata}/profondità desiderata (breve, media, lunga, estesa).
\end{itemize}
A ciò si aggiunge \code{descrizioniI18n}, che replica la matrice per lingua
(\textbf{IT/EN/FR}). Questa struttura è la chiave della \textbf{flessibilità del
modello di annotazione}: lo stesso oggetto si adatta automaticamente al profilo
dell'utente e alla lingua, e aggiungere un livello o una lingua è un'estensione
dei dati, non una modifica del codice. Gli oggetti hanno inoltre un campo
\code{connessi} che definisce un \textbf{grafo} di adiacenze fra opere, usato per
calcolare percorsi (visita una sala alla volta, minimizzando gli spostamenti).

\subsection{Le collezioni MongoDB}
Il dato è partizionato per responsabilità: \code{musei} e \code{musei\_layout}
(struttura e mappa), \code{users} e \code{sessions} (identità), \code{professor\_codes}
e \code{qr\_codes} (autorizzazioni), \code{guided\_visits} (visite di classe),
\code{marketplace\_object\_requests} (richieste d'acquisto),
\code{oggetti\_immagini}/\code{stanze\_immagini} (media). Le route REST riflettono
questa organizzazione: gruppi \code{/musei}, \code{/users}, \code{/guided-visits},
\code{/admin}, \code{/ai}, \code{/qr}.

\subsection{Estendibilità}
Le scelte rendono il sistema facilmente estensibile: un nuovo \emph{device} o
browser è coperto dal layout responsive; un nuovo \emph{modello di dati} (es.
nuovo tipo di oggetto) si aggiunge come campo del documento; un nuovo
\emph{provider IA} o una nuova \emph{lingua} sono configurazioni. La presenza di
percorsi di \emph{fallback} (vedi §\ref{sec:ia}) garantisce che il sistema
degradi in modo controllato anziché rompersi.

% =====================================================================
\section{Comunicazione, autenticazione e sicurezza}

\subsection{Tre livelli di autorizzazione}
\begin{enumerate}[leftmargin=1.5em]
  \item \textbf{Chiave applicativa} (\code{X-API-Key}): ogni richiesta all'API la
    deve portare. Il browser non la conosce: la inietta il BFF (o il proxy Vite
    in sviluppo). Protegge l'API a livello di trasporto.
  \item \textbf{Sessione utente}: dopo il login l'API rilascia un cookie
    \code{muto\_auth} \textbf{HttpOnly} con \code{SameSite=Lax}; il token è
    persistito in \code{sessions} con scadenza. Le password sono salvate come
    hash \textbf{scrypt} con salt, mai in chiaro.
  \item \textbf{Ruoli}: \code{utente}, \code{professore}, \code{admin}. Il ruolo
    professore si ottiene in registrazione tramite un \textbf{codice} (verificato
    come hash SHA-256 lato server); abilita la creazione di visite guidate.
    L'admin accede alle funzioni di gestione del marketplace.
\end{enumerate}
In aggiunta, contenuti specifici si \textbf{sbloccano via QR}: il viewer ha un
componente \code{QrGate} che scansiona il codice fisicamente presente in museo.

\subsection{TLS interno e terminazione esterna}
L'API espone HTTPS anche internamente (il BFF vi si connette in HTTPS, fidandosi
del certificato self-signed). Verso l'esterno, invece, sul cluster di
dipartimento il \textbf{TLS è terminato dal proxy}: il BFF serve HTTP semplice
sulla porta richiesta. Questa distinzione è gestita da configurazione
(\code{BFF\_FORCE\_HTTP}), un esempio di adattabilità all'ambiente di deploy
senza modifiche al codice.

% =====================================================================
\section{Usabilità}

L'attenzione all'utente che \emph{non conosce} il modello applicativo si vede in
diversi punti:
\begin{itemize}[leftmargin=1.3em]
  \item \textbf{Internazionalizzazione completa} IT/EN/FR su tutte le interfacce,
    con la lingua legata al profilo e propagata via evento globale.
  \item \textbf{Adattamento al profilo}: descrizioni e percorsi si calibrano su
    livello e durata scelti, così l'utente vede contenuti adatti senza dover
    capire il modello sottostante.
  \item \textbf{Guida e feedback}: schermate di \emph{briefing} a inizio percorso,
    \emph{toast} non invasivi per esiti delle azioni, stati di caricamento
    (\emph{skeleton}), e ``guardie'' che mostrano messaggi chiari quando manca
    l'autenticazione o il museo.
  \item \textbf{Attività altrui}: nelle visite guidate il professore vede in
    tempo reale lo stato degli studenti (dashboard), e gli studenti seguono la
    navigazione condivisa --- consapevolezza delle azioni altrui.
  \item \textbf{Modalità d'accesso multiple}: chat in linguaggio naturale con
    l'IA (su singolo oggetto o sull'intero museo) e \textbf{comandi vocali}
    offline (riconoscimento con \emph{Vosk}, eseguito interamente nel browser),
    che abbassano la barriera d'uso.
\end{itemize}

% =====================================================================
\section{Sofisticazione grafica}

La presentazione è governata da un \textbf{design system} coerente: un set di
variabili CSS (palette, raggi, tipografia con i font \emph{Cinzel} e
\emph{Raleway}) condiviso fra le SPA, che garantisce uniformità e rende banale un
re-styling globale. Le scelte rilevanti rispetto al rapporto fra dati e
maschere:
\begin{itemize}[leftmargin=1.3em]
  \item \textbf{Mappa vettoriale}: la sede della visita è un SVG generato dal
    layout reale del museo, su cui il viewer disegna percorso, focus e
    spostamenti --- la grafica \emph{è} il dato spaziale, non una decorazione.
  \item \textbf{Differenziazione dei tipi}: oggetti ``normali'' e oggetti di solo
    testo, opere acquistabili e incluse, percorsi standard e percorsi IA
    personalizzati sono distinti visivamente (colore, etichette, badge), così
    che il tipo di dato sia leggibile a colpo d'occhio.
  \item \textbf{Rapporto maschera/dato}: liste con paginazione, gallerie
    immagini con miniature, tabelle delle descrizioni (livello $\times$ durata)
    presentano dati potenzialmente densi in maschere dimensionate al contenuto,
    con \emph{label} comprensibili accanto ai valori formalizzati.
  \item \textbf{Responsive design}: griglie fluide e \emph{breakpoint} adattano
    le maschere a desktop, tablet e mobile.
\end{itemize}

% =====================================================================
\section{Integrazione dell'IA}
\label{sec:ia}

L'IA è usata in modo \emph{assistivo} e sempre con rete di sicurezza:
\begin{itemize}[leftmargin=1.3em]
  \item \textbf{Percorsi personalizzati}: dato un museo e le preferenze
    dell'utente (interessi, livello, durata) e un eventuale \textbf{vincolo in
    linguaggio naturale} (``ho solo mezz'ora'', ``siamo con bambini di 5 e 8
    anni''), il modello seleziona e ordina le opere e genera testi di
    accompagnamento, tradotti nelle tre lingue.
  \item \textbf{Chat contestuale}: domande su un singolo oggetto o sull'intero
    museo, ancorate ai dati realmente presenti per ridurre le allucinazioni.
  \item \textbf{Robustezza}: ogni chiamata IA ha \emph{timeout} e un
    \emph{fallback} deterministico (selezione e descrizioni generate dai dati
    locali). Se il provider è irraggiungibile, l'utente ottiene comunque un
    percorso valido. È un esempio concreto di \emph{flessibilità}: la feature
    degrada, non fallisce.
\end{itemize}

% =====================================================================
\section{Build e deployment}

La pipeline è semplice e riproducibile: \code{npm install} (le dipendenze non
sono versionate) e \code{npm run build} compilano le tre SPA nei rispettivi
\code{dist/}; \code{npm start} avvia il BFF, che a sua volta lancia i servizi
interni (API e SVG server). In produzione sul cluster di dipartimento, il
processo Node ascolta in HTTP sulla porta pubblica e il TLS è gestito a monte;
MongoDB è un container raggiungibile solo dall'interno del cluster. La
configurazione sensibile (chiavi, credenziali) vive in file \code{.env} non
versionati, separando codice e segreti.

\bigskip
\noindent\textbf{In sintesi.} Il progetto privilegia la \emph{tecnologia minima
adeguata} a ogni compito (Alpine dove basta, React dove serve), una netta
separazione di responsabilità fra servizi che comunicano via JSON, un modello
dei dati pensato per adattarsi a profili, lingue e annotazioni diverse, e un
livello di presentazione coerente e consapevole del rapporto fra dati e
maschere. Le soluzioni di compatibilità (HTTP/HTTPS, provider IA, chiave API)
sono risolte per \emph{configurazione}, non forzate nel codice.

\end{document}
